\chapter{Практическая часть}
\label{ch:chap2}

\section{Реализация M-th power carrier recovery}


\begin{lstlisting}
    void synchronizer::coarse_freq_offset(
    const std::vector<std::complex<double>> &samples, rx_cfg &rx_config,
    int Fs) {
    
    std::vector<std::complex<double>> pow_samples(samples.size());

    // pow samples to M-th power
    for (int i = 0; i < samples.size(); ++i) {
        pow_samples[i] = std::pow(samples[i], rx_config.mod_order);
    }

    // calculate FFT
    rx_config.CFO_spectrum = fft(pow_samples, Fs);
    int nfft = rx_config.CFO_spectrum.first.size();

    // find FFT peak frequency
    double max = -__DBL_MAX__;
    int index;

    for (int i = 0; i < nfft; ++i) {
        if (std::abs(rx_config.CFO_spectrum.first[i]) > max) {
            max = std::abs(rx_config.CFO_spectrum.first[i]);
            index = i;
        }
    }

    // calculate frequency offset
    double df = rx_config.CFO_spectrum.second[index] / rx_config.mod_order;
    double dw = -2 * M_PI * df / Fs;

    rx_config.post_cfo_signal.resize(rx_config.rx_samples.size());

    // recovery signal
    for (int i = 0; i < rx_config.rx_samples.size(); ++i) {
        rx_config.post_cfo_signal[i] =
            samples[i] * std::exp(std::complex<double>(0.0, dw * i));
    }

    // optional: calculate FFT after recovery (check result)
    for (int i = 0; i < rx_config.post_cfo_signal.size(); ++i) {
    pow_samples[i] =
        std::pow(rx_config.post_cfo_signal[i], rx_config.mod_order);
    }

    rx_config.post_CFO_spectrum = fft(pow_samples, Fs);
}
\end{lstlisting}

Здесь структура rx\_config - структура для межпоточного взаимодействия GUI и основного потока. FFT - дорогая операция, которая
требует значительной оптимизации, поэтому я использовал библиотеку fftw3. \\

В результате работы функции у нас должно быть 2 спектра: первый спектр, показывающий частотный сдвиг, и второй спектр, показывающий
отсутствие частотного сдвига после восстановления:

\begin{figure}[H]
    \centering
    \includegraphics[width=1.0\textwidth]{bcfo.png}
    \caption{До CFO}
\end{figure}

\begin{figure}[H]
    \centering
    \includegraphics[width=1.0\textwidth]{pcfo.png}
    \caption{После CFO}
\end{figure}

На первом изображении частотный сдвиг примерно 1.8кГц, а на втором уже примерно 0. Это означает, что крупный частотный сдвиг устранен, но
все равно остается небольшой сдвиг + постоянные колебания этого сдвига, которые заставляют фазу крутиться.


\section{Costas Loop}

\begin{lstlisting}
    std::vector<std::complex<double>>
    synchronizer::costas_loop(const std::vector<std::complex<double>> &samples,
                          const int L, const double BnTs, const double Kp,
                          const int mod_order) {

    // define params for loop filter
    const double zeta = std::sqrt(2) / 2;

    double theta = BnTs / (zeta + 0.25 / zeta);
    double denom = 1 + 2 * zeta * theta + theta * theta;

    double K1 = 4 * zeta * theta / (denom * Kp);
    double K2 = 4 * theta * theta / (denom * Kp);

    double phase = 0.0;
    double freq = 0.0;

    double err;

    std::vector<std::complex<double>> out(samples.size());

    auto sign = [](double x) { return (x > 0) - (x < 0); };

    for (size_t i = 0; i < samples.size(); ++i) {
        std::complex<double> rot(std::cos(phase), -std::sin(phase));

        std::complex<double> sample = samples[i];

        // rotate
        out[i] = sample * rot;

        double re = out[i].real();
        double im = out[i].imag();

        //compute error
        if (mod_order == 2) {
        double d = (re + im >= 0) ? 1.0 : -1.0;
        err = d * (im - re);
        } else {
        err = sign(re) * im - sign(im) * re;
        }

        //loop filter
        freq += K2 * err;
        phase += freq + K1 * err;
    }

    return out;
}
\end{lstlisting}

После символьной синхронизации получаем созвездие, которое очень быстро вращается, образовывая при этом окружность. После петли Костаса
вращение должно устраниться, а точки должны скопиться в кучки, образовывая созвездие.

Пример работы на QPSK:

\begin{figure}[H]
    \centering
    \includegraphics[width=1.0\textwidth]{cl.png}
    \caption{После Costas Loop}
\end{figure}

На верхнем изображении на самом деле очень быстрое вращение (скриншот не может это передать), а на нижнем рисунке вращения нет и точки
собрались в QPSK созвездие.

\endinput
